\chapter{Global potentials and positive cohomology classes}
\label{chap:global-potentials}

Let $(X,\omega)$ be a connected compact K\"ahler manifold of dimension $n\geq1$. A positive current has local plurisubharmonic potentials, but their differences need not glue to a function. The obstruction is a finite-dimensional cohomology class. Solving one scalar trace equation identifies that class and produces a global potential. Local compactness and the approximation of the preceding chapter then describe how positivity depends on the class, including classes on the boundary of the pseudo-effective cone.

\section{The trace equation and global potentials}
\label{sec:global-trace-equation}
The K\"ahler form defines the Riemannian metric $g(v,w)=\omega(v,Jw)$ and volume $\mathrm dV=\omega^n/n!$. We use the nonnegative Riemannian Laplacian $L=\mathrm d^*\mathrm d$ on functions. Thus $L=-\sum_j\partial^2/\partial x_j^2$ in Euclidean coordinates; its sign is opposite to the scalar Laplacian used for subharmonic functions. On forms put $\mathcal L=\mathrm d\mathrm d^*+\mathrm d^*\mathrm d$. The Hodge star is characterized by $u\wedge *\overline v=\langle u,v\rangle\mathrm dV$. Stokes' formula gives the adjoint $\mathrm d^*=-*\mathrm d*$ in the even real dimension $2n$.

\begin{lemma}\label{lem:compact-elliptic-foundations}
Distributional forms annihilated by $\mathcal L$ are smooth, and their space is finite dimensional. For every scalar distribution $a$ with $\langle a,1\rangle=0$, there is a unique distribution $f$ with mean zero such that $Lf=a$. If $a$ is smooth, then $f$ is smooth.
\end{lemma}
\begin{proof}
We supply the analytic ingredients. Choose a finite coordinate cover, bundle frames and smooth cutoffs supported inside the charts. For an integer $r\geq0$, define $H^r$ by completing smooth sections in the sum of the chart norms
\[
 \int_{\mathbb R^{2n}}(1+|\xi|^2)^r
       |\widehat{\chi u}(\xi)|^2\,\mathrm d\xi.
\]
\Cref{thm:plancherel} identifies this with the norm of the derivatives of order at most $r$. The chain rule and bounded derivatives on the finitely many compact chart supports show that other such choices give equivalent norms. Set $H^{-r}=(H^r)^*$ using the volume pairing. Fourier inversion and Cauchy--Schwarz give $H^r\subset C^l$ when $r>l+n$: the integral of $|\xi|^{2l}(1+|\xi|^2)^{-r}$ is finite. In particular every distribution, whose order is finite on this compact manifold, belongs to some $H^{-r}$.

The inclusion $H^1\subset L^2$ is compact. In a chart the Fourier tail of an $H^1$-bounded family above frequency $R$ has squared $L^2$ norm at most $C/R^2$. Its frequency-truncated inverses are uniformly bounded and equicontinuous on a fixed coordinate box, by Cauchy--Schwarz on $|\xi|\leq R$. Arzel\`a--Ascoli gives compactness there. The localized original functions vanish outside the box; approximation on the box and the Fourier tail bound therefore give $L^2$ compactness. A finite chart argument proves the assertion.

For completeness we also prove the elliptic regularity needed here. The principal symbol of $\mathrm d$ is exterior multiplication by $\mathrm i\xi$, and that of its adjoint is its Hermitian adjoint. Exterior multiplication and contraction satisfy
\[
 \iota_\xi(\xi\wedge u)+\xi\wedge\iota_\xi u=|\xi|_g^2u.
\]
Consequently $\mathcal L$ has scalar positive principal part on each coefficient. Freezing the coefficients on a sufficiently small chart and using Plancherel bounds the gradient by the quadratic form of $\mathcal L$. The change in the principal coefficients has arbitrarily small supremum and is absorbed into this gradient bound. Derivatives of the coefficients and cutoffs produce terms bounded by $C\|u\|_{L^2}\|u\|_{H^1}$; the inequality $ab\leq\varepsilon a^2+b^2/(4\varepsilon)$ absorbs these as well. Summing a finite partition gives
\begin{equation}\label{eq:compact-garding}
 \|u\|_{H^1}^2\leq C\bigl(\|\mathrm du\|_{L^2}^2+
       \|\mathrm d^*u\|_{L^2}^2+\|u\|_{L^2}^2\bigr).
\end{equation}
Thus, for a fixed sufficiently large $c>0$, the quadratic form of $P=\mathcal L+c$ is an equivalent Hilbert norm on $H^1$. Riesz representation solves $Pv=b$ uniquely in $H^1$ for $b\in H^{-1}$.

Here is the regularity of this weak solution, without assuming a smooth solution in advance. In a chart write its principal divergence form as $-\partial_i(a^{ij}\partial_j v)$, with $a^{ij}$ scalar and uniformly positive. Write $\delta_h v(x)=(v(x+h e_k)-v(x))/h$. In the weak equation use the test $-\delta_{-h}(\chi^2\delta_hv)$, first with a smooth approximation to $v$. Discrete integration by parts and ellipticity give
\[
 \|\chi\nabla\delta_hv\|_{L^2}^2
 \leq C\bigl(\|b\|_{L^2}+\|v\|_{H^1}\bigr)
             \bigl(\|\chi\nabla\delta_hv\|_{L^2}+\|v\|_{H^1}\bigr).
\]
Indeed $\|\delta_hv\|_{L^2}\leq\|\partial_kv\|_{L^2}$, coefficient differences are bounded by their first derivatives, and the derivatives of $\chi$ and the lower-order terms are bounded. The right side containing $b$ is bounded by its $L^2$ norm times the $L^2$ norm of the test; the latter is bounded by the second factor displayed above, using the same difference-quotient inequality. Approximation in $H^1$ makes this test legitimate. Absorption gives a bound independent of $h$. Weak $L^2$ limits of $\delta_h\partial_jv$ are the distributional derivatives, so $v\in H^2$ on the smaller chart. Differentiating the equation $r$ times gives the same estimate for those derivatives: its commutators involve derivatives of $v$ of order at most $r+1$. Induction, nested cutoffs and the finite cover show
\[
 \|P^{-1}b\|_{H^{r+2}}\leq C_r\|b\|_{H^r},\qquad r\geq0.
\]
The weak estimate covers $r=-1$. Self-adjointness extends $P^{-1}$ by duality to $H^r\to H^{r+2}$ for every negative integer $r$ as well. The two inverse identities extend by density and duality. If $\mathcal Lu=0$ distributionally, then $u=cP^{-1}u$; starting from $H^{-r}$ and iterating improves regularity until all positive Sobolev orders are reached. The Fourier embedding makes $u$ smooth. For smooth harmonic forms, \eqref{eq:compact-garding} bounds the $H^1$ norm by the $L^2$ norm. Their closed $L^2$ unit ball is compact by the preceding compactness argument, so the harmonic space is finite dimensional: an infinite orthonormal sequence would contradict compactness.

On mean-zero scalar functions, $\|v\|_{L^2}\leq C\|\mathrm dv\|_{L^2}$. Otherwise a sequence of norm one with gradients tending to zero has an $L^2$ limit of norm one by compactness; its distributional gradient vanishes, hence it is constant on the connected manifold, contradicting its zero mean. Riesz representation with the gradient norm therefore solves $Lf=a$ for smooth mean-zero $a$. Testing by $f$ and the regularity just proved show that its inverse $G$ is continuous from $C^\infty$ to $C^\infty$, self-adjoint, and satisfies $LG=GL=1-\Pi$, where $\Pi$ is averaging onto constants. Extend $G$ to distributions by transposition. These identities still hold, proving existence for every mean-zero distribution. A distribution in the kernel of $L$ is smooth by the form regularity, and integration of $|\mathrm df|^2$ makes it constant. This gives uniqueness.
\end{proof}

Let $\Lambda$ denote contraction by $\omega$, so that $\Lambda\omega=n$. In a unitary holomorphic coframe a real $(1,1)$-form is a Hermitian matrix; $\Lambda$ is its matrix trace.

\begin{lemma}\label{lem:kahler-trace-star}
For a scalar distribution $f$ and a real $(1,1)$-current $T$,
\begin{equation}\label{eq:kahler-trace-star}
 \Lambda\ddc f=-\frac{Lf}{4\pi},\qquad
 *T=(\Lambda T)\frac{\omega^{n-1}}{(n-1)!}
          -T\wedge\frac{\omega^{n-2}}{(n-2)!}\quad(n\geq2).
\end{equation}
If $\mathrm dT=0$ and $\Lambda T$ is constant, then $\mathcal LT=0$.
\end{lemma}
\begin{proof}
At a point choose holomorphic coordinates with $\omega=\frac{\mathrm i}{2}\sum_j\mathrm dz_j\wedge\mathrm d\bar z_j$ and vanishing first metric derivatives. Such coordinates exist because $\mathrm d\omega=0$ gives $\partial_jg_{k\bar l}=\partial_kg_{j\bar l}$: a quadratic holomorphic change removes those symmetric first derivatives. There $L=-4\sum_j\partial_j\partial_{\bar j}$ and contraction of $\ddc f$ is $\pi^{-1}\sum_j f_{j\bar j}$. This proves the first identity. Diagonalize the Hermitian matrix of $T$ by a unitary transformation. For its $j$th diagonal basis term, the star is the product of the other $n-1$ basis terms; expansion of the two wedges in the second identity gives exactly that product. Linearity proves the identity. These are identities of differential operators and smooth bundle maps, hence hold for distributions. Closedness of $\omega$ now gives $\mathrm d*T=0$ when the trace is constant. Thus $\mathrm dT=\mathrm d^*T=0$. In dimension one a $(1,1)$-current is $(\Lambda T)\omega$, so the last assertion holds there too.
\end{proof}

\Needspace{12\baselineskip}
\begin{theorem}[The global $dd^c$ lemma]\label{thm:global-ddc-lemma}
Every closed real $(1,1)$-current $T$ has a unique decomposition
\begin{equation}\label{eq:global-current-decomposition}
 T=H+\ddc f,
\end{equation}
where $H$ is a smooth harmonic real $(1,1)$-form and $f$ is a scalar distribution of mean zero. If $T$ is smooth, $f$ is smooth. If $T$ is positive, $f$ has a canonical quasi-plurisubharmonic representative. Moreover $T$ is exact as a current if and only if $H=0$.
\end{theorem}
\begin{proof}
Put $a=\int_XT\wedge\omega^{n-1}/\int_X\omega^n$. The identity $T\wedge\omega^{n-1}=(\Lambda T)\omega^n/n$ shows that $\Lambda T-na$ has mean zero. Solve $Lf=-4\pi(\Lambda T-na)$ by \cref{lem:compact-elliptic-foundations}. Then $U=T-a\omega-\ddc f$ is closed with zero trace. \Cref{lem:kahler-trace-star} makes $U$ harmonic, and \cref{lem:compact-elliptic-foundations} makes it smooth. Set $H=a\omega+U$.

An exact harmonic form $B$ is zero: if $B=\mathrm dS$, its squared norm equals $\langle\mathrm dS,*B\rangle=0$ since $\mathrm d*B=0$. This argument works for a distributional $S$. The difference between two candidate $H$'s is exact, so they coincide; then $\ddc(f-f')=0$ and the trace equation makes $f-f'$ constant. The normalization gives uniqueness. It also proves the exactness criterion. Smoothness of $T$ gives smoothness of $f$ by the scalar equation.

If $T\geq0$, on a ball write $T=\ddc u$ with $u$ plurisubharmonic and $H=\ddc h$ with $h$ smooth, using \cref{thm:local-positive-current-potential,lem:smooth-potential-on-ball}. Then $\ddc(f-u+h)=0$. The local scalar Laplacian of this distribution vanishes, so Weyl's lemma makes it a smooth pluriharmonic function. Thus $f$ is locally a plurisubharmonic function minus a smooth function. The canonical local representatives agree: their almost-everywhere equality is equality of plurisubharmonic representatives after adding a common smooth potential. They define the stated global representative.
\end{proof}

\begin{definition}\label{def:real-one-one-cohomology}
The real $(1,1)$-cohomology space\index{cohomology!real $(1,1)$ classes} is
\[
 H^{1,1}(X,\mathbb R)=
 \frac{\{\text{closed smooth real $(1,1)$-forms}\}}
      {\{\ddc f:f\in C^\infty(X,\mathbb R)\}}.
\]
Write $\{T\}=\{H\}$ for the class of a closed real $(1,1)$-current as in \eqref{eq:global-current-decomposition}.
\end{definition}
Real de Rham cohomology in degree two is the quotient of closed smooth real two-forms by forms $\mathrm d\beta$ with $\beta$ a smooth real one-form. \Cref{thm:global-ddc-lemma} identifies the space just defined with the harmonic real $(1,1)$-forms and with the subspace of de Rham cohomology represented by $(1,1)$-forms. In particular it is finite dimensional. We give it that vector-space topology. Choosing a harmonic basis provides a linear smooth representative $\alpha\mapsto\theta_\alpha$; there is a norm and a constant $C$ such that
\begin{equation}\label{eq:small-class-small-form}
 -C\|\alpha\|\omega\leq\theta_\alpha\leq C\|\alpha\|\omega.
\end{equation}
The class of a current is unchanged by adding $\ddc$ of a distribution. Harmonic coefficients are its pairings against smooth harmonic forms, so the class map on closed currents is weakly continuous.

If $\theta$ is a smooth representative of $\{T\}$ and $T\geq0$, the same theorem gives $T=\theta+\ddc\varphi$ with $\varphi\in\PSH(X,\theta)$, unique up to a constant. Changing $\theta$ by $\ddc h$ changes $\varphi$ by $-h$ up to a constant. Thus the polar locus, Lelong numbers and singularity comparisons of $T$ are independent of the smooth representative.

\section{Normalized potentials and minimal singularities}
\label{sec:global-minimal-potentials}
Fix a closed smooth real $(1,1)$-form $\theta$. Local compactness controls its potentials once constants have been fixed on the connected manifold.

\begin{proposition}\label{prop:normalized-global-psh-compactness}
For $a\leq b$ the set of $\varphi\in\PSH(X,\theta)$ with $a\leq\sup_X\varphi\leq b$ is compact in $L^1(X)$. The supremum is continuous for this topology. In particular the normalized family $\sup_X\varphi=0$ has a uniform $L^1$ bound. These statements remain valid for sequences of backgrounds converging smoothly, with the limit potential taken relative to the limiting background.
\end{proposition}
\begin{proof}
On a finite cover add smooth local potentials for $\theta$, bounded on the smaller charts. The resulting plurisubharmonic functions are uniformly bounded above. Take points of maximum and pass to a convergent subsequence. On a chart containing their limit the functions cannot converge locally uniformly to $-\infty$, since their values at those points are bounded below. \Cref{thm:psh-compactness-hartogs} therefore gives a noncollapsed $L^1$ limit there. On an overlapping chart, collapse would contradict that limit on the nonempty open intersection. Connectedness and a finite chain of overlaps propagate noncollapse over the finite cover. Diagonal selection gives a global $L^1$ limit, and distributional positivity supplies its canonical $\theta$-psh representative.

If $\varphi_j\to\varphi$ in $L^1$, local ball averages give uniform upper bounds. For $x_j\to x$, the continuous-majorant form of Hartogs' lemma gives $\limsup_j\varphi_j(x_j)\leq\varphi(x)$. Indeed choose a small neighborhood of $x$ on which $\varphi$ is below the constant $\varphi(x)+\varepsilon$, and add the smooth local potential to that majorant before applying the lemma; when $\varphi(x)=-\infty$, use any finite constant instead. Applied at maximum points this bounds the upper limit of the suprema. For a subsequence on which the suprema tend to their lower limit, almost-everywhere convergence gives $\varphi\leq\liminf_j\sup_X\varphi_j$ almost everywhere and hence everywhere. The two inequalities prove continuity and the closedness of the stated interval constraint. Sequential compactness is compactness in the metric space $L^1$; a continuous norm is bounded on its compact normalized subset.

For smoothly convergent backgrounds use one common lower bound $-C\omega$ and fixed local potentials for $C\omega$ to obtain compactness. Distributional convergence identifies the limiting background. The same Hartogs argument, after adding these fixed local potentials, proves the assertion about suprema.
\end{proof}

\begin{corollary}\label{cor:nonpluripolar-normalization}
If a compact set $K\subset X$ is non-pluripolar, then $\sup_X\varphi-\sup_K\varphi$ is uniformly bounded for $\varphi\in\PSH(X,\theta)$. Consequently a lower and upper bound for $\sup_K\varphi$ gives an $L^1$-relatively compact family.
\end{corollary}
\begin{proof}
Otherwise choose normalized $\varphi_j\leq0$ with $\sup_K\varphi_j\leq-2^j$. The partial sums $\sum_{j=1}^N2^{-j}\varphi_j$ have a common lower Hessian bound $-C\omega$, are decreasing, and have uniformly bounded $L^1$ norms by \cref{prop:normalized-global-psh-compactness}. Their limit is a quasi-psh function, by the local decreasing-limit theorem, and equals $-\infty$ on $K$. This contradicts non-pluripolarity. The bound just proved converts the bounds on $K$ into bounds on $X$.
\end{proof}

\begin{corollary}\label{cor:global-decreasing-nets}
A decreasing net $\varphi_i\in\PSH(X,\theta)$ whose suprema are bounded below has a $\theta$-psh infimum $\varphi$, converges to it in $L^1$, and satisfies $\sup_X\varphi_i\to\sup_X\varphi$.
\end{corollary}
\begin{proof}
On a tail the suprema are bounded above as well. Compactness supplies an $L^1$ cluster subnet. Its limit $\psi$ is below every $\varphi_i$ almost everywhere and therefore pointwise, by cofinality and the canonical representative. Thus the pointwise infimum is not identically $-\infty$. Apply the local decreasing-net theorem on the finite cover to obtain its representative and $L^1$ convergence; continuity of the supremum gives the last assertion. The subnet here is a cofinal order-preserving Willard subnet, as in the global conventions.
\end{proof}

\begin{definition}\label{def:global-minimal-potential}
If $\PSH(X,\theta)\ne\varnothing$, put
\[
 V_\theta=\bigl(\sup\{\varphi\in\PSH(X,\theta):\varphi\leq0\}\bigr)^*.
\]
A potential $\psi\in\PSH(X,\theta)$ has \emph{minimal singularities}\index{minimal singularities} if every $\theta$-psh potential is more singular than it. The symbol $^*$ denotes upper semicontinuous regularization.
\end{definition}
\Cref{thm:psh-upper-envelopes}, applied after adding smooth local potentials, gives $V_\theta\in\PSH(X,\theta)$, $V_\theta\leq0$ and $\sup_XV_\theta=0$. Every $\varphi\in\PSH(X,\theta)$ satisfies $\varphi\leq V_\theta+\sup_X\varphi$, so $V_\theta$ has minimal singularities. Any two minimal potentials satisfy both comparison inequalities with constants; their difference is bounded on their common finite locus. The definition is unchanged by replacing $\theta$ within its class, since the change of potential is smooth and bounded.

\begin{proposition}\label{prop:smooth-target-psh-modification}
Let $\pi:Y\to X$ be a proper bimeromorphic morphism of connected compact K\"ahler manifolds. Pullback is a bijection $\PSH(X,\theta)\to\PSH(Y,\pi^*\theta)$ and is continuous for $L^1$ convergence. For every prime divisor $E$ on a fixed such model, the Lelong numbers $\nu(\varphi,E)$ are uniformly bounded for $\varphi\in\PSH(X,\theta)$.
\end{proposition}
\begin{proof}
There is an analytic $A\subset X$ of codimension at least two over whose complement $\pi$ is an isomorphism, by \cref{prop:smooth-bimeromorphic-divisor-locus}. A potential on $Y$ descends there. It is locally bounded above near $A$: properness confines the inverse image of a small compact neighborhood to a compact subset of $Y$, on which the potential is bounded above. Adding a local potential for $\theta$, bounded removal across $A$ gives a plurisubharmonic extension. Its pullback equals the original potential off the exceptional set and hence everywhere by equality of canonical representatives. Pullback in the other direction is defined by holomorphic pullback and cannot be identically $-\infty$, since it is finite almost everywhere on the isomorphism locus. This also proves injectivity.

For normalized convergent potentials, their pullbacks have the same supremum zero and a fixed background, hence are relatively compact by \cref{prop:normalized-global-psh-compactness}. On the isomorphism locus every cluster limit is the pullback of the given limit, by change of variables on compact subsets. Canonical representatives identify it everywhere, proving continuity; the continuity of the suprema handles unnormalized sequences. The generic divisorial Lelong number is an upper semicontinuous function on this compact family: on the fixed model it is the infimum over points of the upper semicontinuous point Lelong numbers from \cref{prop:lelong-calculus,prop:generic-divisorial-lelong}. It is finite there, so it is bounded above. Constants do not change it.
\end{proof}

\begin{corollary}\label{cor:global-psh-balayage}
Let $B\Subset B'$ be coordinate balls and let $\varphi\in\PSH(X,\theta)$ be bounded near $\overline{B'}$. There is a canonical $\theta$-psh replacement $\psi\geq\varphi$, equal to $\varphi$ on $X\setminus B$, bounded near $\overline B$, and satisfying $(\theta+\ddc\psi)^n=0$ on $B$. It is the largest such replacement; among continuous replacements it is unique.
\end{corollary}
\begin{proof}
Write $\theta=\ddc h$ on $B'$ and apply \cref{cor:homogeneous-psh-balayage} to $\varphi+h$. Subtract $h$ and retain $\varphi$ outside $B$. Every replacement has upper boundary limits at most the specified boundary values by upper semicontinuity, so the boundary-envelope characterization proves the largest-member assertion. If a replacement is continuous, its boundary values are continuous and uniqueness is \cref{thm:homogeneous-ball-dirichlet}. The product is the bounded Bedford--Taylor product in the local potential.
\end{proof}

\section{Positivity cones and divisorial decomposition}
\label{sec:positive-cohomology-cones}
For a class $\alpha$ fix a smooth representative $\theta$. The global potential theorem makes positivity of its currents equivalent to inequalities for $\theta+\ddc\varphi$.

\begin{definition}\label{def:cohomology-positivity-cones}
A class $\alpha\in H^{1,1}(X,\mathbb R)$ is \emph{pseudo-effective}\index{cohomology!pseudo-effective class} if it contains a closed positive current, \emph{K\"ahler}\index{cohomology!K\"ahler class} if it contains a K\"ahler form, and \emph{nef}\index{cohomology!nef class} if for every $\varepsilon>0$ it has a smooth representative at least $-\varepsilon\omega$. It is \emph{big}\index{cohomology!big class} if it contains a closed current at least $\delta\omega$ for some $\delta>0$; such a current is called a \emph{K\"ahler current}\index{K\"ahler current}. The corresponding cones are denoted $\mathcal E$, $\mathcal K$, $\mathcal N$ and $\mathcal B$. Zero is included in the pseudo-effective and nef cones.
\end{definition}
These definitions are independent of $\omega$: any two smooth positive metrics on the compact manifold bound each other by fixed multiples. Pseudo-effectivity is equivalent to $\PSH(X,\theta)\ne\varnothing$.

\begin{proposition}\label{prop:cohomology-cone-topology}
The pseudo-effective cone is closed and convex, the K\"ahler cone is open and convex, $\mathcal N=\overline{\mathcal K}\subset\mathcal E$, $\mathcal K$ is the interior of $\mathcal N$, and $\mathcal B$ is the interior of $\mathcal E$.
\end{proposition}
\begin{proof}
Convexity follows by addition and positive scaling of representatives. A positive current $T$ has total trace mass
\[
 \int_XT\wedge\omega^{n-1}
       =\int_X\theta_{\{T\}}\wedge\omega^{n-1},
\]
since the distributional $\ddc$ term pairs to zero against the closed form. For a convergent sequence of classes this mass is bounded. The positive-current coefficient bounds and weak measure compactness from \cref{thm:positive-current-order-zero} give a weakly convergent subsequence of their positive representatives. The limit is positive and closed; weak continuity of the harmonic pairings identifies its class. This proves closedness of $\mathcal E$.

The small-representative bound \eqref{eq:small-class-small-form} preserves strict positivity under a sufficiently small class change, proving openness of $\mathcal K$. A nef representative at least $-\varepsilon\omega$ becomes positive after adding $2\varepsilon\omega$; hence its class is in $\overline{\mathcal K}$. Conversely perturb representatives of convergent K\"ahler classes by the small harmonic difference in \eqref{eq:small-class-small-form}; this gives the nef inequalities. Closedness of $\mathcal E$ gives $\mathcal N\subset\mathcal E$. Adding any positive multiple of $\{\omega\}$ to a nef class gives a K\"ahler class, by taking the negative allowance smaller than that multiple. An interior nef class remains nef after subtracting a sufficiently small multiple, and is therefore K\"ahler. If $T\in\alpha$ dominates $\delta\omega$, the same small-representative bound puts a neighborhood of $\alpha$ in $\mathcal E$. Conversely an interior class satisfies $\alpha-\delta\{\omega\}\in\mathcal E$ for sufficiently small $\delta>0$. Adding $\delta\omega$ to a positive representative proves bigness.
\end{proof}

\begin{proposition}\label{prop:bounded-potential-nef}
A class containing a positive current with locally bounded potentials is nef. A class containing a K\"ahler current with locally bounded potentials is K\"ahler. Every big class contains a K\"ahler current with gentle rational analytic singularities.
\end{proposition}
\begin{proof}
Write the given current $\theta+\ddc\varphi$. For bounded $\varphi$, every $\mathcal I(m\varphi)$ is the unit ideal. The cutoff-sum approximations of \cref{lem:logarithmic-cutoff-sum} are therefore smooth, with lower bound $-c/m$ for $\theta+\ddc G_m$. This gives nefness. For a current at least $\delta\omega$, use the approximation with background $\theta-\delta\omega$, giving $\theta+\ddc G_m\geq(\delta-c/m)\omega$. This is a K\"ahler form for large $m$. Without boundedness the same inequality supplies a K\"ahler current with the gentle analytic type proved in that lemma.
\end{proof}

\begin{corollary}\label{cor:strict-minimal-current-kahler}
A positive current with minimal singularities can be a K\"ahler current only if its class is K\"ahler. A K\"ahler class has bounded minimal potentials.
\end{corollary}
\begin{proof}
Write the minimal current as $T=\theta+\ddc\varphi\geq\delta\omega$ and normalize $\varphi\leq0$. For sufficiently small $t>0$, the current $(1-t)T+t\theta=\theta+\ddc((1-t)\varphi)$ remains positive. Minimality gives $(1-t)\varphi\leq\varphi+C$, so $\varphi\geq-C/t$ at its finite points. This holds almost everywhere, and the canonical local representative extends the lower bound to every point. Thus $\varphi$ is bounded. Apply \cref{prop:bounded-potential-nef} in its strict regime. Conversely a K\"ahler form in the class has a smooth potential relative to $\theta$; minimality bounds the normalized minimal potential below by that smooth competitor.
\end{proof}

For a prime divisor $D$, the generic Lelong number $\nu(T,D)$ is the generic number of any local potential. Smooth potentials contribute zero, so this is independent of the background.

\Needspace{10\baselineskip}
\begin{theorem}[Siu's divisorial decomposition]\label{thm:siu-divisorial-decomposition}
Every closed positive $(1,1)$-current has a unique decomposition
\begin{equation}\label{eq:siu-divisorial-decomposition}
 T=R+\sum_D\nu(T,D)[D],
\end{equation}
with $R\geq0$ closed and $\nu(R,D)=0$ for every prime divisor. The positive series has at most countably many nonzero terms and converges as a current with measure coefficients. Thus $T-\nu(T,D)[D]\geq0$.
\end{theorem}
\begin{proof}
Put $c=\nu(T,D)$. Near a smooth point of $D$ choose $D=\{z_1=0\}$ and a local plurisubharmonic potential $u$. \Cref{prop:generic-lelong-uniform-growth} gives $u\leq c\log|z_1|^2+C$. Off $D$, the function $u-c\log|z_1|^2$ is plurisubharmonic and bounded above near $D$. Bounded analytic removal and Poincar\'e--Lelong therefore give a positive extension of $T-c[D]$ near $D_{\mathrm{reg}}$. On $D_{\mathrm{sing}}$, of codimension at least two, the measure coefficients of $[D]$ vanish, since they are obtained by integration on $D_{\mathrm{reg}}$. The restriction of the coefficient measures of $T-c[D]$ to that set is thus the restriction of $T$, and is positive. Testing the coefficient measures on its complement and on this set separately proves positivity everywhere. In dimension one the same local argument removes the atom of mass $c$ at the point divisor.

By \cref{thm:siu-lelong-level-sets}, each $\{\nu(T,x)\geq1/j\}$ is a proper analytic subset. Its compactness and \cref{cor:analytic-global-components} give finitely many divisorial components. Every divisor with $\nu(T,D)>0$ is a component of one of these sets, so there are countably many. Subtract a finite list successively. At a generic point of a new divisor, the earlier divisors are absent, so its generic coefficient has not changed; the first part applies again. Thus every finite subtraction is positive. Their masses are bounded by the mass of $T$; the positive coefficient bounds give convergence of the increasing series and its closed positive residual $R$.

For a fixed divisor $D$, write the series as its $D$ term and the other positive terms. Additivity and positivity of Lelong numbers give $\nu(R,D)+\nu(T,D)\leq\nu(T,D)$, so $\nu(R,D)=0$. In any decomposition of the stated form, the series $S_D$ of the other divisors has zero coefficient mass on $D$: each intersection of two distinct prime divisors has zero integration measure, and a countable sum preserves this property. The first part gives $S_D\geq\nu(S_D,D)[D]$, forcing $\nu(S_D,D)=0$. Finite additivity applied to $R+b_D[D]+S_D$ therefore identifies its coefficient $b_D$ with $\nu(T,D)$. This proves uniqueness.
\end{proof}

\begin{definition}\label{def:modified-nef-class}
A class $\alpha$ is \emph{modified nef}\index{cohomology!modified nef class} if for every $\varepsilon>0$ there is a closed current $T_\varepsilon\in\alpha$ with $T_\varepsilon\geq-\varepsilon\omega$ and $\nu(T_\varepsilon+\varepsilon\omega,D)=0$ for every prime divisor $D$. Write $\mathcal{MN}$ for this cone. A positive current with zero generic divisorial Lelong numbers is \emph{non-divisorial}\index{current!non-divisorial}.
\end{definition}
The cone $\mathcal{MN}$ is convex, contains $\mathcal N$, and is contained in $\mathcal E$. For the last inclusion use positive representatives of $\alpha+\varepsilon\{\omega\}$ and closedness of $\mathcal E$. It is closed as well: for $\alpha_j\to\alpha$ choose $j$ with a smooth representative of $\alpha-\alpha_j$ bounded by $\varepsilon\omega/2$, and add it to a representative of $\alpha_j$ bounded below by $-\varepsilon\omega/2$ with all divisorial numbers zero. Smooth additions do not change those numbers. Commensurability of metrics proves independence from $\omega$. Applying \cref{lem:logarithmic-cutoff-sum} to $T_\varepsilon+\varepsilon\omega$ shows that the representatives can have gentle analytic singularities with pole sets of codimension at least two: at a generic point of each divisor the bound $\operatorname{ord}\mathcal I(mu)\leq m\nu(u)=0$ makes the ideal a unit, so the approximation has no divisorial poles. For the prescribed allowance use the original representative with allowance $\varepsilon/2$ and choose $m$ with $c/m<\varepsilon/2$.

\begin{corollary}\label{cor:divisorial-logarithmic-residual}
If $\theta+\ddc\varphi\geq0$ and $\varphi$ has logarithmic singularities along an effective rational divisor $D$, with bounded residual local potentials after subtracting its logarithms, then $\{\theta\}-\{D\}$ is nef. If the original current is K\"ahler, that residual class is K\"ahler.
\end{corollary}
\begin{proof}
The proof of \cref{thm:siu-divisorial-decomposition} identifies the residual current $T-[D]$ as positive with bounded local potentials. In the strict case first subtract $\delta\omega$ from $T$; its divisorial numbers are unchanged, so the residual still dominates $\delta\omega$. Apply \cref{prop:bounded-potential-nef}.
\end{proof}

\section{Minimal multiplicities and positivity loci}
\label{sec:non-nef-non-kahler}
For $\alpha\in\mathcal E$ and $\theta\in\alpha$, the class $\alpha+\varepsilon\{\omega\}$ is big for every $\varepsilon>0$. Its minimal potential measures the positivity remaining after a small allowance in curvature.

\begin{definition}\label{def:minimal-class-multiplicities}
For $x\in X$ and a prime divisor $D$ put
\[
 \nu(\alpha,x)=\sup_{\varepsilon>0}\nu(V_{\theta+\varepsilon\omega},x),
 \qquad
 \nu(\alpha,D)=\sup_{\varepsilon>0}\nu(V_{\theta+\varepsilon\omega},D).
\]
These are the \emph{minimal multiplicities}\index{minimal multiplicity}. Define the \emph{non-nef locus}\index{non-nef locus} by
\[
 E_{\mathrm{nn}}(\alpha)=\{x:\nu(\alpha,x)>0\}.
\]
For a big class define its \emph{non-K\"ahler locus}\index{non-K\"ahler locus} by
\[
 E_{\mathrm{nK}}(\alpha)=\bigcap_T\{\varphi_T=-\infty\},
\]
where $T=\theta+\ddc\varphi_T$ ranges over its K\"ahler currents with rational analytic singularities. Set $E_{\mathrm{nK}}(\alpha)=X$ for a class which is not big, and $E_{\mathrm{nn}}(\alpha)=X$ for a class which is not pseudo-effective.
\end{definition}
The minimal potentials decrease as $\varepsilon\downarrow0$, so their Lelong numbers increase. The suprema can therefore be taken along $\varepsilon=1/j$. Changing the smooth representative shifts potentials by a bounded smooth function, and changing $\omega$ sandwiches their singularities between the allowances for fixed multiples of $\varepsilon$. Thus the definitions are intrinsic. Every positive current in $\alpha$ is a competitor at every allowance; hence both multiplicities are finite and bounded by its corresponding Lelong numbers. Outside a countable union of proper analytic subsets of $D$, all the generic numbers of the sequence equal its point numbers. Consequently
\[
 \nu(\alpha,D)=\inf_{x\in D}\nu(\alpha,x),
\]
and this infimum is attained at very general points of $D$.

\begin{proposition}\label{prop:big-class-minimal-multiplicity}
If $\alpha$ is big, then $\nu(\alpha,x)=\nu(V_\theta,x)$ and $\nu(\alpha,D)=\nu(V_\theta,D)$. For an arbitrary pseudo-effective class only the inequalities with $\leq$ hold. Moreover $\alpha$ is modified nef if and only if $\nu(\alpha,D)=0$ for every prime divisor $D$.
\end{proposition}
\begin{proof}
The positive current $T_0=\theta+\ddc V_\theta$ is a competitor with every negative allowance, giving the inequalities. In the big case choose $S\in\alpha$ with $S\geq\delta\omega$. Write $T_\varepsilon=\theta+\ddc V_{\theta+\varepsilon\omega}\geq-\varepsilon\omega$ and $t=\varepsilon/(\delta+\varepsilon)$. Then $(1-t)T_\varepsilon+tS$ is a positive current in $\alpha$. Minimality gives
\[
 \nu(T_0,x)\leq(1-t)\nu(T_\varepsilon,x)+t\nu(S,x).
\]
Let $\varepsilon\downarrow0$; finiteness of $\nu(S,x)$ proves equality. The identical inequality with generic numbers proves the divisorial assertion. If $\alpha$ is modified nef, its zero-divisorial representatives are competitors for the minimal potentials with every allowance, so all the generic minimal multiplicities are zero. Conversely, if those multiplicities vanish, the currents $T_\varepsilon$ themselves satisfy the definition of modified nefness.
\end{proof}

\begin{proposition}\label{prop:minimal-multiplicity-class-continuity}
For a fixed point or prime divisor, minimal multiplicity is homogeneous, convex and lower semicontinuous on $\mathcal E$, and continuous on $\mathcal B$.
\end{proposition}
\begin{proof}
Scaling a current and its curvature allowance proves homogeneity. Adding currents minimal with allowance $\varepsilon$ in two classes gives a competitor with allowance $2\varepsilon$ in their sum. Minimality and additivity of Lelong numbers therefore give subadditivity; passage to the decreasing allowances gives convexity. If $\alpha_j\to\alpha$, fix $\varepsilon>0$ and add to a minimal current in $\alpha_j$ with allowance $\varepsilon$ a smooth representative of $\alpha-\alpha_j$ bounded in modulus by $\varepsilon\omega$. It becomes a competitor for the minimal current in $\alpha$ with allowance $2\varepsilon$. Thus that minimal number is at most $\liminf_j\nu(\alpha_j,x)$. Taking the supremum over $\varepsilon$ gives lower semicontinuity; the same argument applies to $D$.

A finite convex function is continuous on an open convex subset of a finite-dimensional vector space. Here is the elementary local reason. In a small cube inside that subset, convexity bounds the function above by the maximum of its finitely many vertex values. Reflection in its center gives a lower bound on a smaller cube. Extend the segment between two points of a still smaller cube by a fixed distance inside the first cube; the convex slope inequality bounds their value difference by a constant times their distance. Apply this to the open cone $\mathcal B$.
\end{proof}

\begin{theorem}\label{thm:positivity-loci}
For $\alpha\in\mathcal E$, the non-nef locus is a countable union of proper analytic subsets and is empty precisely when $\alpha$ is nef. If $\alpha$ is big, its non-K\"ahler locus is the pole set of one rational analytic K\"ahler current, is proper analytic, contains its non-nef locus, and is empty precisely when $\alpha$ is K\"ahler. For every class,
\begin{equation}\label{eq:non-kahler-non-nef-allowance}
 E_{\mathrm{nK}}(\alpha)=
       \bigcap_{\varepsilon>0}E_{\mathrm{nn}}(\alpha-\varepsilon\{\omega\}).
\end{equation}
\end{theorem}
\begin{proof}
The non-nef locus is the union, over positive integers $j,k$, of the analytic sets $\{\nu(V_{\theta+\omega/j},x)\geq1/k\}$. Each is proper by \cref{thm:siu-lelong-level-sets}. If $\alpha$ is nef, then $\alpha+\varepsilon\{\omega\}$ is K\"ahler and has a smooth competitor. Its minimal potential is bounded, so all these point numbers vanish. Conversely suppose they vanish. The order bound $\operatorname{ord}_x\mathcal I(mV_{\theta+\varepsilon\omega})\leq m\nu(V_{\theta+\varepsilon\omega},x)=0$ makes the ideal a unit at every point. The smooth cutoff-sum approximation has background $\theta+\varepsilon\omega$ and curvature loss $c/m$. Hence $\alpha$ has smooth representatives at least $-(\varepsilon+c/m)\omega$, proving nefness.

For a big class analytic K\"ahler currents exist by \cref{prop:bounded-potential-nef}. Their pole sets are closed analytic. Second countability selects countably many complements with the same union as all the complements: take the countable basis elements contained in one of the open complements and choose one complement for each such element. Finite intersections of the selected pole sets form a descending sequence of analytic sets. \Cref{lem:strong-analytic-noetherian}, applied to the increasing sums of their coherent defining ideals on the compact manifold, makes it stationary. Thus their intersection is already a finite intersection. Choose the corresponding currents, a common strict lower bound $\delta\omega$, and their potentials relative to $\theta$. Their finite maximum is still $\theta-\delta\omega$-psh and has rational analytic singularities by \cref{prop:rational-analytic-operations}; its pole set is exactly the intersection. This gives one current realizing $E_{\mathrm{nK}}(\alpha)$, and its pole set is proper since its potential is finite almost everywhere. Minimality and \cref{prop:big-class-minimal-multiplicity} show that every positive Lelong point of the minimal current belongs to that pole set.

If the realized pole set is empty, its analytic potential is bounded. \Cref{prop:bounded-potential-nef} then gives a K\"ahler form in the class. The converse uses that form itself. Finally, if $x$ is outside $E_{\mathrm{nK}}(\alpha)$, choose an analytic K\"ahler current at least $\delta\omega$ whose potential is bounded near $x$. For $0<\varepsilon<\delta$ the current minus $\varepsilon\omega$ is positive in $\alpha-\varepsilon\{\omega\}$ and has zero Lelong number at $x$. It bounds that class's minimal multiplicity by zero.

Conversely, suppose $x\notin E_{\mathrm{nn}}(\alpha-\varepsilon\{\omega\})$. That class is pseudo-effective by the definition outside $\mathcal E$. For $0<\eta<\varepsilon/2$, put $u=V_{\theta-(\varepsilon-\eta)\omega}$. Its Lelong number at $x$ is zero. The ideal $\mathcal I(mu)$ is a unit near $x$, and its global cutoff-sum approximation $G_m$ has
\[
 \theta+\ddc G_m\geq(\varepsilon-\eta-c/m)\omega.
\]
For large $m$ this is an analytic K\"ahler current whose potential is smooth near $x$. Thus $x$ is outside $E_{\mathrm{nK}}(\alpha)$; this also proves that the class is big whenever the right-hand intersection in \eqref{eq:non-kahler-non-nef-allowance} omits a point. The convention for non-big classes is therefore consistent with the formula.
\end{proof}

\begin{proposition}\label{prop:positivity-under-modifications}
For a proper bimeromorphic morphism $\pi:Y\to X$ of connected compact K\"ahler manifolds, pullback preserves nef classes and pushforward preserves modified nef classes. Pushforward of a non-divisorial positive current is non-divisorial.
\end{proposition}
\begin{proof}
Fix a K\"ahler form $\omega_Y$. Smoothness and compactness give $\pi^*\omega\leq C\omega_Y$. Pullback of the nef representatives with their lower bounds therefore proves the first assertion. Pushforward of a closed current is closed by testing against pullbacks; it sends distributional exact currents to exact currents, so induces the stated class map. At a generic point of each prime divisor $D\subset X$, $\pi$ is an isomorphism by \cref{prop:smooth-bimeromorphic-divisor-locus}. Its strict transform $D'$ has $\nu(\pi_*S,D)=\nu(S,D')$, proving the assertion about non-divisorial positive currents.

If $\beta$ is modified nef, $\beta+t\{\omega_Y\}$ is big and modified nef for $t>0$. Its positive minimal current is non-divisorial by \cref{prop:big-class-minimal-multiplicity}; its pushforward is thus a positive non-divisorial current, whose class is modified nef. As $t\downarrow0$ these classes tend to $\pi_*\beta$. Closedness of $\mathcal{MN}$ proves the claim.
\end{proof}

\begin{example}\label{ex:nef-boundary-minimal-current}
Boundary multiplicities must be defined with the negative curvature allowances. The always valid comparison is $\nu(\alpha,x)\leq\nu(V_\theta,x)$; the proof of equality in \cref{prop:big-class-minimal-multiplicity} uses a strict current. In contrast, $E_{\mathrm{nK}}(\alpha)=X$ for every non-big class, including the zero class, although $E_{\mathrm{nn}}(0)=\varnothing$. Thus the two loci distinguish different failures of positivity at the boundary.
\end{example}

\section{Curvature and metrics defined by sections}
\label{sec:global-line-bundle-metrics}
Let $L\to X$ be a holomorphic line bundle and fix a smooth quadratic metric $h_0$ as in \cref{def:quadratic-line-bundle-metric}. Its curvature is denoted $c_1(L,h_0)$, and its cohomology class $c_1(L)$ is independent of the metric: if $h_1=h_0\mathrm e^{-f}$ then $c_1(L,h_1)=c_1(L,h_0)+\ddc f$. Tensor products add the curvatures and pullback pulls them back. These statements follow in local frames by addition and composition of their logarithmic squared norms.

\begin{definition}\label{def:psh-quadratic-line-metric}
A \emph{plurisubharmonic metric}\index{metric!plurisubharmonic} on $L$ is the quadratic metric $h=h_0\mathrm e^{-\varphi}$ for $\varphi\in\PSH(X,c_1(L,h_0))$. Its value on the zero vector is zero, including at poles; on a nonzero vector over a pole it is $+\infty$. Its Chern current is
\[
 c_1(L,h)=c_1(L,h_0)+\ddc\varphi.
\]
\end{definition}
This is a definition using the canonical pointwise potential. The distributional curvature inequality for an arbitrary measurable weight only specifies its almost-everywhere class. Changing $h_0$ makes the corresponding smooth change in $\varphi$, so the definition and its pointwise metric do not depend on $h_0$. Thus $L$ has such a metric precisely when $c_1(L)$ is pseudo-effective. The tensor product of two psh metrics is defined by adding their local weights; a pole in either factor makes the nonzero-vector norm infinite. Holomorphic pullback is a psh metric when the pulled-back weights are not identically $-\infty$ on a component. Restriction to a submanifold uses that same condition.

\begin{proposition}\label{prop:line-bundle-poincare-lelong}
For a nonzero holomorphic section $s$ of $L$ and a smooth metric $h_0$,
\begin{equation}\label{eq:line-bundle-poincare-lelong}
 c_1(L,h_0)+\ddc\log|s|_{h_0}^2=[Z(s)].
\end{equation}
\end{proposition}
\begin{proof}
In a local frame $s=f e$, the left side is $\ddc\log|f|^2$, because $\log|s|_{h_0}^2=\log|f|^2+\log|e|_{h_0}^2$. Apply \cref{thm:poincare-lelong} and glue the local zero divisors. In particular the integration current represents $c_1(L)$ with coefficient one.
\end{proof}

\Needspace{8\baselineskip}
\begin{example}\label{ex:projective-fubini-study-curvature}
For $\mathcal O_{\mathbb P^N}(1)$, dualizing the Euclidean norm on the tautological line gives the Fubini--Study metric. On the chart $[1:z_1:\cdots:z_N]$, the coordinate section $X_0$ has squared norm $(1+|z|^2)^{-1}$ and curvature
\[
 \omega_{\mathrm{FS}}=\ddc\log(1+|z|^2).
\]
Its Hermitian Hessian has matrix $((1+|z|^2)\delta_{jk}-\bar z_jz_k)/(1+|z|^2)^2$, which is positive definite by Cauchy--Schwarz. On $\mathbb P^1$, polar integration gives $\int\omega_{\mathrm{FS}}=1$. The potential $\log(|z|^2/(1+|z|^2))$, with value zero at infinity, defines a singular psh metric whose Chern current is $[0]$, by \eqref{eq:line-bundle-poincare-lelong}.
\end{example}

\Needspace{14\baselineskip}
\begin{definition}\label{def:section-fubini-study-metrics}
A \emph{Fubini--Study metric}\index{metric!Fubini--Study} on $L$ is $h_0\mathrm e^{-\varphi}$ with
\[
 \varphi=\frac1m\max_{1\leq i\leq N}
          \bigl(\log|s_i|_{h_0^m}^2+\lambda_i\bigr),
 \qquad m\geq1,\quad \lambda_i\in\mathbb Q,
\]
where the $s_i$ are holomorphic sections of $L^m$ and at least one is nonzero. Denote their set by $\FS(L)$. The \emph{generalized Fubini--Study metrics}\index{metric!generalized Fubini--Study} form the smallest subset of psh metrics containing $\FS(L)$ and closed under minima of metrics and rational convex combinations of weights. A convex combination of weights corresponds to the geometric combination of quadratic metrics, with value zero on the zero vector and the usual infinite value at a pole of either positive-weight factor.
\end{definition}

\begin{proposition}\label{prop:fubini-study-finite-operations}
These metrics are independent of $h_0$, have rational analytic singularities, and are closed under the two specified operations. Thus, for these finite rational operations, generalized Fubini--Study metrics coincide with $\FS(L)$. The set may be empty even when the definition of a psh metric is applicable.
\end{proposition}
\begin{proof}
In a frame $e$ write $s_i=f_i e^m$. The local psh weight of the metric is $m^{-1}\max_i(\log|f_i|^2+\lambda_i)$, independent of the reference metric. It differs by a bounded function from $m^{-1}\log\sum_i\mathrm e^{\lambda_i}|f_i|^2$. This proves positivity, rational analytic singularities and independence, including the common zero set.

For two presentations of degrees $m,l$, raise their sections to powers to obtain a common degree $M$; the associated constants are multiplied by the same integer powers and remain rational. The maximum of the two potentials is represented by the union of these section lists, and corresponds to the minimum of the metrics. For a rational $t=a/b\in(0,1)$, the weighted sum of the two maxima is the maximum over all pairs of their weighted sums. The sections $s_i^a r_j^{b-a}$ have degree $Mb$, and their constants $a\lambda_i+(b-a)\mu_j$ are rational. They give exactly that convex combination. Induction proves closure under any finite number of operations and hence the equality of the generated sets. If no positive tensor power has a nonzero section there is no such presentation; the definition makes no nonemptiness claim in that case.
\end{proof}

For a point, the cohomology and all four positivity cones are $\{0\}$, positivity is vacuous, both loci are empty, and normalized potentials are zero. Metrics are quadratic norms. On a disconnected compact manifold apply the constructions separately to its finitely many components, imposing normalization and noncollapse on each.

\section{Compact interpolation and global Bergman kernels}
\label{sec:compact-point-bergman}
The local extension estimate extends from a point on a compact manifold once the curvature dominates the cost of a global logarithmic pole. We first establish the bundle estimate used to correct the cutoff. All norms in this section use the volume $\mathrm dV=\omega^n/n!$. The canonical bundle $K_X=\Lambda^{n,0}T^*X$ is the holomorphic line bundle of top holomorphic covectors; its metric and the dual metric on $K_X^{-1}$ are induced by $\omega$. We write $H^0(X,L)$ for the vector space of global holomorphic sections of $L$.

\begin{lemma}\label{lem:analytic-complement-cutoffs}
Let $A$ be the closed proper pole set of a quasi-plurisubharmonic function $\psi$ on $X$, smooth on $X\setminus A$, with $\psi\to-\infty$ at every point of $A$. There is a K\"ahler form $\sigma\geq\omega/2$ on $X\setminus A$ with smooth compactly supported cutoffs $\chi_R\to1$, $0\leq\chi_R\leq1$, and $|\mathrm d\chi_R|_\sigma\leq C/R$. In particular this applies to finite unions of gentle analytic pole sets and point poles.
\end{lemma}
\begin{proof}
Choose $b$ with $\ddc\psi\geq-b\omega$ and subtract a constant so that $\psi<-2b-2$. For $\rho=-\log(-\psi)$ the chain rule gives
\[
 \ddc\rho=-\frac1\psi\ddc\psi+
       \frac{\mathrm i}{2\pi\psi^2}\partial\psi\wedge\bar\partial\psi.
\]
Thus $\sigma=\omega+\ddc\rho\geq\omega/2+
(\mathrm i/2\pi)\partial\rho\wedge\bar\partial\rho$. The last term bounds $|\mathrm d\rho|_\sigma$ by a numerical constant: in a unitary coframe it is a rank-one Hermitian form whose sole eigenvalue is the squared norm of $\partial\rho$, with the indicated factor. The function $-\rho$ tends to $+\infty$ at $A$ and has compact sublevels in $X\setminus A$. Take a smooth function $\chi$ equal to one on $(-\infty,1]$ and zero on $[2,\infty)$, and put $\chi_R=\chi(-\rho/R)$. These are the required cutoffs. The same gradient bound shows that a path approaching $A$ has infinite $\sigma$-length, since $-\rho$ is unbounded along it. The form is therefore complete as well. For the last assertion add the finitely many gentle potentials and logarithmic point potentials; their sum has precisely the union as its pole set, the stated lower bound, and the required divergence.
\end{proof}

\begin{lemma}\label{lem:kahler-bundle-l2-estimate}
Let $M$ be an open subset of $X$, let $F\to M$ be a holomorphic line bundle with a smooth quadratic metric $h$, and suppose $c_1(F,h)\geq a\omega$ for $a>0$. Suppose also that $M$ has a K\"ahler metric $\sigma\geq\omega/2$ and cutoffs as in \cref{lem:analytic-complement-cutoffs}. For a distributionally closed $F$-valued $(n,1)$-form $g$ whose squared norm is integrable with respect to $\omega,h$, there is an $F$-valued $(n,0)$-form $u$ with $\bar\partial u=g$ and
\begin{equation}\label{eq:compact-bundle-l2-estimate}
 \int_M|u|_{\omega,h}^2\,\mathrm dV
       \leq\frac1{2\pi a}\int_M|g|_{\omega,h}^2\,\mathrm dV.
\end{equation}
\end{lemma}
\Needspace{10\baselineskip}
\begin{proof}
We give the differential and domain arguments. In a holomorphic frame of squared norm $\mathrm e^{-p}$ the Chern operators are $D'=\partial-\partial p\wedge$ and $D''=\bar\partial$. They act on bundle-valued forms, and $D'D''+D''D'=\partial\bar\partial p\wedge$. Write $\Lambda_\sigma$ for contraction by $\sigma$. In holomorphic normal coordinates and a frame with $\partial p=0$ at the chosen point, exterior multiplication by the unitary covectors and their contractions satisfy
\[
 \iota_j\epsilon_k+\epsilon_k\iota_j=\delta_{jk},\qquad
 D''{}^*=-\mathrm i[\Lambda_\sigma,D'],\qquad
 D'{}^*=\mathrm i[\Lambda_\sigma,D''].
\]
The barred covectors have the same anticommutation rule. To check the two operator identities, expand $D'$ and $D''$ as covector multiplication followed by the corresponding coefficient derivative; moving a contraction past that covector leaves its Kronecker term, which is the negative coefficient derivative in the formal adjoint. First metric derivatives vanish in these coordinates. The identities hold at every point, hence globally. Substitution, cancellation of the two mixed compositions and the curvature identity give
\[
 \begin{gathered}
 \Delta''=\Delta'+[2\pi c_1(F,h),\Lambda_\sigma],
 \\
 \Delta''=D''D''{}^*+D''{}^*D'',\qquad
 \Delta'=D'D'{}^*+D'{}^*D'.
 \end{gathered}
\]
Here the commutator with a two-form means wedge multiplication followed by contraction minus the reverse composition. Set $B_\sigma=[2\pi c_1(F,h),\Lambda_\sigma]$. For a compactly supported $(n,1)$-form $v$, integration yields
\begin{equation}\label{eq:kahler-bundle-bochner}
 \|D''{}^*v\|_\sigma^2+\|D''v\|_\sigma^2
 =\|D'{}^*v\|_\sigma^2+
       \int_M\langle B_\sigma v,v\rangle_\sigma\,\mathrm dV_\sigma.
\end{equation}
The term $D'v$ is zero by its holomorphic degree. Diagonalizing the curvature in a $\sigma$-unitary coframe shows that $B_\sigma$ acts on the coefficient of its $j$th barred covector by the $j$th curvature eigenvalue. It is positive.

Use the maximal distributional domains for $T=D'':L^2_{(n,0)}\to L^2_{(n,1)}$ and $S=D'':L^2_{(n,1)}\to L^2_{(n,2)}$, with the $\sigma,h$ norms. They are closed and densely defined, and $ST=0$. Compactly supported smooth forms are dense in $\operatorname{Dom}T^*\cap\operatorname{Dom}S$ in graph norm. The cutoff product rules are proved from the adjoint identity exactly as in \cref{lem:three-weight-density}; their error norms are at most $C\|v\|_\sigma/R$, so the cutoffs give compact supports. Local smoothing also works for the variable coefficients here. Indeed, for $P=\sum a_l(x)\partial_l+b(x)$ its commutator with convolution has first-order part
\[
 \sum_l\int(a_l(x)-a_l(y))\partial_l\rho_t(x-y)v(y)\,\mathrm dy
   +\sum_l\int(\partial_l a_l)(y)\rho_t(x-y)v(y)\,\mathrm dy.
\]
On a fixed compact set the mean-value bound for $a_l$ and scaling of $\rho_t$ make these operators uniformly bounded on $L^2$. For smooth $v$ their sum tends to zero; density of smooth coefficients in $L^2$ proves the same for every $v$. The zeroth-order commutator tends to zero by ordinary local $L^2$ convergence. Apply this simultaneously to $T^*$ and $S$, then a finite partition of unity. This proves graph density without imposing an unstated boundary condition.

Pointwise duality bounds $|\langle g,v\rangle_\sigma|^2$ by the product of
$\int\langle B_\sigma^{-1}g,g\rangle_\sigma\,\mathrm dV_\sigma$ and the last integral in \eqref{eq:kahler-bundle-bochner}. The first integral is independent of the auxiliary metric $\sigma$. To see this, write the $(n,1)$-form as a fixed top covector times a barred one-form. Its top-degree squared norm times volume is metric independent, while inversion of $B_\sigma$ cancels the inverse-metric factor on the barred coefficient. The resulting dual form is the inverse of the curvature form $2\pi c_1(F,h)$ itself. Since that form is at least $2\pi a\omega$, this integral is at most the right side of \eqref{eq:compact-bundle-l2-estimate}. Also $g$ belongs to $L^2_{(n,1)}(\sigma,h)$: the top-degree density is unchanged and the barred covector norm is at most twice its $\omega$ norm. Graph approximation passes the Bochner inequality to the operator domains: local convergence and Fatou's lemma control its nonnegative curvature integral. The inequality and \cref{thm:hilbert-complex-solvability} now give $Tu=g$ with this bound. The $(n,0)$ norm times volume is likewise independent of the K\"ahler metric, giving exactly \eqref{eq:compact-bundle-l2-estimate}.
\end{proof}
\index{Bochner--Kodaira identity!K\"ahler line bundle}

\begin{lemma}\label{lem:l2-holomorphic-analytic-removal}
A holomorphic function on $U\setminus A$, where $A$ is a proper analytic subset of a complex manifold $U$, extends holomorphically to $U$ if it belongs to $L^2_{\mathrm{loc}}(U)$ after arbitrary definition on $A$. The same holds for sections of a holomorphic line bundle.
\end{lemma}
\begin{proof}
Near a point of $A$ choose a nonzero holomorphic function $b$ vanishing on $A$. Choose a coordinate basis on which the first nonzero homogeneous term of $b$ is nonzero in each basis direction. Weierstrass preparation then makes the zero set of $b$ finite on each coordinate-line fiber of a sufficiently small polydisc; exceptional parameter fibers form a null set. By Fubini the given function is $L^2$ on almost every such line. Its Laurent expansion at each puncture has no negative coefficient: circle orthogonality would otherwise give a divergent radial integral $\int_0^r t^{1-2j}\,\mathrm dt$ for some $j\geq1$. Thus the slice extends holomorphically and has zero distributional barred derivative on that line. Fubini against compactly supported tests shows that every barred coordinate derivative of the original $L^2$ function is zero distributionally on the polydisc. The distributional holomorphicity consequence of Weyl's lemma in \cref{sec:weighted-dbar-estimate} gives the extension. It agrees with the original on $U\setminus A$ and glues by uniqueness. Local holomorphic frames prove the bundle assertion.
\end{proof}

\begin{theorem}\label{thm:compact-point-extension}
There are constants $N,C>0$, depending only on $(X,\omega)$, with the following property. Let $L$ be any holomorphic line bundle, $h_0$ any smooth quadratic metric and $h=h_0\mathrm e^{-\varphi}$ a psh metric with $c_1(L,h)\geq N\omega$. For $x\in X$ with $\varphi(x)>-\infty$ and $v\in L_x$ there is $s\in H^0(X,L)$ such that
\begin{equation}\label{eq:compact-point-extension}
 s(x)=v,\qquad
 \int_X|s|_h^2\,\mathrm dV\leq C|v|_h^2.
\end{equation}
The constants are independent of $L,h_0,\varphi,x,v$.
\end{theorem}
\begin{proof}
Fix finitely many coordinate balls with nested smaller balls covering $X$. Translated balls around $x$ inside these charts can have a common positive radius. A cutoff equal to one on such an inner ball and zero outside a larger one, multiplied by $n\log|z-z(x)|^2$, gives a global function $\psi_x\leq0$. Choose the scale in the logarithm so it is negative on the chart. Its curvature is positive on the inner ball; on the cutoff annulus the logarithm and its first two derivatives are uniformly bounded. Thus
\[
 \ddc\psi_x\geq-b_0\omega,
 \qquad \psi_x=n\log|z-z(x)|^2+O(1)\text{ near }x,
\]
with $b_0\geq0$ independent of $x$. Choose $b_1\geq0$ with $c_1(K_X^{-1},h_{K^{-1},\omega})\geq-b_1\omega$, where the canonical metric is induced by $\omega$, and take $N=b_0+b_1+2$.

In a frame of $L$ on the selected ball, the local weight of $h$ is plurisubharmonic. \Cref{cor:ot-point-value-bound} supplies a holomorphic local section $f$ with $f(x)=v$ and
\[
 \int_{B}|f|_h^2\,\mathrm dV\leq C_0|v|_h^2.
\]
The volume comparison is uniform on the finite chart cover. Choose another cutoff $\chi$ supported in $B$ and equal to one on a smaller ball around $x$. Put $G=\chi f$ and $g=\bar\partial G=f\bar\partial\chi$. The support of $g$ stays in a fixed annulus, where $\psi_x$ and the cutoff derivatives have uniform bounds.

Apply \cref{thm:gentle-quasi-equisingular-approximation} with background $c_1(L,h_0)-N\omega$. It gives gentle potentials $\varphi_j\downarrow\varphi$, with analytic pole sets $A_j$ and $c_1(L,h_0)+\ddc\varphi_j\geq(N-\varepsilon_j)\omega$, $\varepsilon_j\to0$. Since $\varphi_j(x)\geq\varphi(x)>-\infty$, none of the $A_j$ contains $x$. Set
\[
 F=L\otimes K_X^{-1},\qquad
 h_{F,j}=h_0\mathrm e^{-\varphi_j-\psi_x}\otimes h_{K^{-1},\omega}.
\]
For late $j$ this metric is smooth on $M_j=X\setminus(A_j\cup\{x\})$ and has curvature at least $\omega$. The cutoffs of \cref{lem:analytic-complement-cutoffs} are available there. The canonical identification of $L$-valued $(0,q)$-forms with $F$-valued $(n,q)$-forms is an isometry for $\omega$, because the top covector and its dual cancel. Apply \cref{lem:kahler-bundle-l2-estimate} to $g$ under this identification. Its weighted norm is bounded uniformly:
\[
 \int_{M_j}|g|_{h_0}^2\mathrm e^{-\varphi_j-\psi_x}\,\mathrm dV
 \leq C_1\int_B|f|_h^2\,\mathrm dV
 \leq C_1C_0|v|_h^2.
\]
Here $\varphi_j\geq\varphi$ is essential. We obtain a correction $u_j$ with $\bar\partial u_j=g$ on $M_j$ and the same uniform bound for
$\int|u_j|_{h_0}^2\mathrm e^{-\varphi_j-\psi_x}$.

On $M_j$, $s_j=G-u_j$ is holomorphic. The weights are bounded below relative to $h_0$, since $\varphi_j$ is bounded above and $\psi_x\leq0$. Thus $s_j$ is locally unweighted $L^2$ across $A_j\cup\{x\}$, and \cref{lem:l2-holomorphic-analytic-removal} extends it to a global section. Near $x$, $g=0$ and $A_j$ is absent; hence $u_j$ is holomorphic there after removal. Its integrability with the factor $|z-z(x)|^{-2n}$ forces $u_j(x)=0$, exactly as in \cref{prop:logarithmic-point-interpolation}. Consequently $s_j(x)=v$. Since $\psi_x\leq0$, the correction estimate also bounds its norm without that pole, and
\[
 \int_X|s_j|_{h_0}^2\mathrm e^{-\varphi_j}\,\mathrm dV
 \leq C_2|v|_h^2.
\]
The constant uses only the fixed geometric cutoffs and the local point estimate.

The uniform upper bound $\varphi_j\leq\varphi_1$ makes the sections bounded in unweighted $L^2(h_0)$. Local submean inequalities and Cauchy estimates, followed by Arzel\`a--Ascoli on the finite cover, give a subsequence converging uniformly to a global holomorphic section $s$. It has $s(x)=v$. Fatou's lemma and $\varphi_j\downarrow\varphi$ give \eqref{eq:compact-point-extension}. The case $v=0$ uses the zero section.
\end{proof}
\index{Ohsawa--Takegoshi theorem!compact point extension}

\Needspace{22\baselineskip}
\begin{proposition}\label{prop:compact-analytic-bergman-approximation}
Let $L,T$ be holomorphic line bundles with smooth quadratic metrics $h_0,h_T$, and put $\theta=c_1(L,h_0)$ and $\theta_T=c_1(T,h_T)$. Suppose $\varphi$ has rational analytic singularities and $\theta+\ddc\varphi\geq\delta\omega$ for $\delta>0$. Define
\[
 \mathcal H_k=H^0\bigl(X,L^k\otimes T\otimes\mathcal I(k\varphi)\bigr),
 \qquad \|s\|_k^2=\int_X|s|_{h_0^k\otimes h_T}^2\mathrm e^{-k\varphi}\,\mathrm dV.
\]
For all sufficiently large $k$, this is a nonzero finite-dimensional Hilbert space, and its Bergman potential is
\begin{equation}\label{eq:compact-bergman-potential}
 \varphi_k=\frac1k\log\sup_{\|s\|_k\leq1}|s|_{h_0^k\otimes h_T}^2.
\end{equation}
It has gentle rational analytic singularities and belongs to $\PSH(X,\theta+k^{-1}\theta_T)$. There are $C>0$ and an increasing sequence $\alpha_k\in(0,1)$ tending to one such that
\begin{equation}\label{eq:compact-bergman-type-bounds}
 \varphi_k\geq\varphi-\frac{\log C}{k},\qquad
 \varphi\preceq\varphi_k\preceq\alpha_k\varphi.
\end{equation}
Here $u\preceq v$ means $u\leq v+O(1)$, as in \cref{sec:analytic-approximation-comparison}.
\end{proposition}
\begin{proof}
Choose $b_T$ with $\theta_T\geq-b_T\omega$. For $k\delta-b_T\geq N$, \cref{thm:compact-point-extension} applies to $L^k\otimes T$ with metric $(h_0^k\otimes h_T)\mathrm e^{-k\varphi}$. It supplies, at each finite point of $\varphi$, a section taking any chosen unit reference-metric value and with squared norm at most $C\mathrm e^{-k\varphi(x)}$. Rescaling proves the first bound in \eqref{eq:compact-bergman-type-bounds}; at poles the bound is automatic. At least one finite point exists, so $\mathcal H_k\ne0$.

The space of all holomorphic sections of a fixed line bundle on $X$ is finite dimensional. Its unweighted $L^2$ unit ball is compact by the same finite-chart submean, Cauchy and Arzel\`a--Ascoli argument just used; an infinite orthonormal sequence would contradict this compactness. Thus $\mathcal H_k$ is finite dimensional. Weighted integrability is exactly the displayed multiplier-ideal condition, and its norm is a positive definite inner product. For an orthonormal basis $s_1,\ldots,s_{d_k}$, Hilbert projection gives the supremum in \eqref{eq:compact-bergman-potential} as $\sum_i|s_i|^2$. In a frame the holomorphic logarithmic tuple therefore gives positivity with background $\theta+k^{-1}\theta_T$, and analytic coefficient $1/k$. Its exponential $\mathrm e^{k\varphi_k}$ is smooth globally. On a log resolution of the coherent base ideal, divide by its common divisor monomial. The resulting holomorphic tuple has no common zero, so its sum of squared norms is smooth and strictly positive. This proves gentleness with the same quadratic-metric convention.

For the upper type bound choose a finite cover by local analytic presentations of $\varphi$ and log resolutions over smaller neighborhoods, as in \cref{lem:analytic-weighted-integrability-growth}. Let $a_E>0$ be their finitely many positive logarithmic coefficients and $A_E$ the corresponding discrepancies. Choose $B>\max_E A_E/a_E$; if there are no such components choose any $B>0$. The proof of \cref{lem:analytic-weighted-integrability-growth} only requires $k\eta a_E\geq A_E$ on this finite list. For $\eta=B/k$ that holds for every component. Its factoring and properness argument therefore gives, for $k>B$ and every local coefficient of every $s_i$,
\[
 \log|s_i|_{h_0^k\otimes h_T}^2
       \leq(k-B)\varphi+C_{i,k}.
\]
Smooth frame factors only change the constants. A finite chart cover and the finite basis give $\varphi_k\leq(1-B/k)\varphi+C_k$. Choose $\alpha_k=1-B/k$ for the stated large integers. It is positive, increases to one, and proves the remaining comparisons. If $\varphi$ is bounded the same conclusion holds because the local ideal is the unit ideal and the right side is bounded. No bounded-growth ideal or division theorem is required.
\end{proof}
\index{Bergman kernel!compact analytic approximation}

\begin{example}\label{ex:projective-line-bergman-kernel}
On $\mathbb P^1$ take $L=\mathcal O(1)$ with its Fubini--Study metric, the trivial twist, $\varphi=0$ and volume $\omega_{\mathrm{FS}}$. A section of $L^k$ is a polynomial of degree at most $k$ on the chart $[1:z]$: holomorphicity of $\zeta^k f(1/\zeta)$ at infinity and Cauchy's estimates exclude all higher coefficients. Rotational integration makes the monomials orthogonal, and
\[
 \|z^j\|_k^2=\int_0^\infty\frac{t^j}{(1+t)^{k+2}}\,\mathrm dt
       =\frac{j!(k-j)!}{(k+1)!},\qquad 0\leq j\leq k.
\]
The substitution $u=t/(1+t)$ followed by integration by parts gives the last equality. Thus the orthonormal homogeneous sections are
$\sqrt{(k+1)\binom{k}{j}}X_0^{k-j}X_1^j$. Their squared norms sum to
\[
 \frac{k+1}{(1+|z|^2)^k}\sum_{j=0}^k\binom{k}{j}|z|^{2j}=k+1.
\]
Consequently $\varphi_k=k^{-1}\log(k+1)$ converges uniformly to zero. The calculation checks both the quadratic metric normalization and the volume $\int_{\mathbb P^1}\omega_{\mathrm{FS}}=1$.
\end{example}
